\subsection*{الف}
بنا بر تعریف کانولوشن پیوسته:
\begin{flalign*}
    (G_\sigma * f)(x,y) &= \iint_{\mathbb{R}^2} G_\sigma(x-u,\, y-v)\, f(u,v)\, \dd u\, \dd v &&
\end{flalign*}
با مشتق‌گیری نسبت به $x$ و انتقال عملگر مشتق به زیر انتگرال (مجاز است، چون $f$ مستقل از $x$ و $G_\sigma$ نسبت به $x$ به‌قدر کافی هموار و انتگرال‌پذیر است):
\begin{flalign*}
    \frac{\partial}{\partial x}(G_\sigma * f)(x,y)
    &= \frac{\partial}{\partial x} \iint G_\sigma(x-u,\, y-v)\, f(u,v)\, \dd u\, \dd v \\
    &= \iint \frac{\partial}{\partial x} G_\sigma(x-u,\, y-v)\, f(u,v)\, \dd u\, \dd v \\
    &= \iint \left(\frac{\partial G_\sigma}{\partial x}\right)(x-u,\, y-v)\, f(u,v)\, \dd u\, \dd v \\
    &= \left(\frac{\partial G_\sigma}{\partial x} * f\right)(x,y) &&
\end{flalign*}
انتقال مشتق به زیر انتگرال با
\lr{Leibniz integral rule}
و قاعدهٔ زنجیری
$\partial_x G_\sigma(x-u,\cdot) = (\partial_x G_\sigma)(x-u,\cdot)$
توجیه می‌شود؛ هموار و سریعاً میرا بودن گاوسی شرط جابه‌جایی را تأمین می‌کند.

\subsection*{ب}
مشتق‌گیری عملگری است که نوفهٔ فرکانس‌بالا را تقویت می‌کند؛ در حوزهٔ فرکانس ضریب آن $\propto |\omega|$ است:
\begin{flalign*}
    \mathcal{F}\{\partial_x f\}(\omega) = j\omega_x\, \hat f(\omega) &&
\end{flalign*}
از این رو مشتق‌گیری مستقیم روی تصویر خام، نوفهٔ موجود را به‌شدت بزرگ‌نمایی می‌کند و لبه‌های واقعی در میان پاسخ‌های کاذب گم می‌شوند. گاوسی یک
\lr{Low-pass filter}
است و حاصل‌ضرب طیف آن با $|\omega|$ یک پاسخ
\lr{Band-pass}
می‌سازد که نوفهٔ فرکانس‌بالا را پیش از مشتق‌گیری سرکوب می‌کند. افزون بر این، مشتق روی شبکهٔ گسسته
\lr{ill-posed}
است و تنها پس از هموارسازی پایدار و معنادار می‌شود؛ پارامتر $\sigma$ نیز مقیاس لبه‌های آشکارشده را کنترل می‌کند.

\subsection*{ج}
عملگر
\lr{Laplacian}
(و
\lr{LoG})
در شرایط زیر بر گرادیان مرتبهٔ اول برتری دارد:
\begin{itemize}
    \item ایزوتروپ بودن: $\nabla^2$ نسبت به دوران ناورداست و به جهت لبه وابسته نیست، در حالی‌که آشکارساز گرادیانی نیازمند ترکیب مؤلفه‌های جهتی است.
    \item مکان‌یابی دقیق لبه از طریق
    \lr{Zero-crossing}؛
    عبور از صفر در $\nabla^2 (G_\sigma * f)=0$ مرز را با ضخامت یک پیکسل و بسته (پیوسته) مشخص می‌کند، حال آن‌که بیشینهٔ گرادیان لبه‌ای ضخیم و نیازمند آستانه‌گذاری می‌دهد.
    \item تشخیص هم‌زمان لبه‌ها در همهٔ جهت‌ها با یک کانولوشن واحد و به‌صرفه از نظر محاسباتی.
    \item مناسب برای آشکارسازی نقاط، خطوط نازک و بافت‌های ریز که تغییر مرتبهٔ دوم در آن‌ها بارزتر است.
\end{itemize}
کاربرد بهینه آنجاست که تصویر پس از هموارسازیِ گاوسی نوفهٔ کنترل‌شده دارد و لبه‌ها به‌صورت گذارهای ملایم شدت ظاهر می‌شوند؛ در این حالت
\lr{LoG}
با انتخاب مناسب $\sigma$ مرزهای بسته و دقیق استخراج می‌کند.